\section{Heuristic functions}\label{sec:heuristics}

The heuristic is the only application-specific part of A*. This section defines
the conditions under which A* is correct, presents the distance heuristics of
our experiments, and shows how a heuristic guides or misleads the search.

\subsection{Requirements}\label{sec:h-requirements}

A good heuristic $h$ meets three requirements.
\begin{enumerate}[label=(\roman*)]
    \item \textbf{It is cheap to compute}, because A* computes it for every
    vertex it reaches.
  \item \textbf{It never overestimates}: $h(v) \le h^{*}(v)$ for every $v$. This
    makes A* return a shortest path (Section~\ref{sec:correctness}).
  \item \textbf{It is close to $h^{*}$}, because larger values leave fewer
    vertices for A* to expand (Section~\ref{sec:efficiency}).
\end{enumerate}
The best heuristic is therefore as large as possible without exceeding the
true cost.

\subsection{Admissible heuristics}\label{sec:admissible}

Requirement~(ii) is called admissibility.

\begin{definition}[Admissible heuristic]\label{def:admissible}
A heuristic $h$ is \emph{admissible} if $0 \le h(v) \le h^{*}(v)$ for every
vertex $v$.
\end{definition}

An admissible heuristic is optimistic. The heuristic of the running example
is admissible (Table~\ref{tab:story-h}); Section~\ref{sec:overestimate} gives
one that is not. At the destination an admissible heuristic is zero, since
$0 \le h(t) \le h^{*}(t) = 0$.

\subsection{Consistent heuristics}\label{sec:consistent}

Algorithm~\ref{alg:astar} needs a stronger condition.

\begin{definition}[Consistent heuristic]\label{def:consistent}
A heuristic $h$ is \emph{consistent} if $h(t) = 0$ and
\begin{equation}\label{eq:consistent}
  h(u) \le w(u,v) + h(v)
\end{equation}
for every edge $(u,v)$.
\end{definition}

In words, one step lowers the estimate by at most the cost of that step, a
triangle inequality (Figure~\ref{fig:consistency}). On undirected graphs it must
hold in both directions of every edge. The heuristic of the running example is
consistent; for instance, $h(\cty{N}) = 5 = w(\cty{N}, \cty{B}) + h(\cty{B})$.

\begin{figure}[htbp]
  \centering
  \begin{tikzpicture}[x=1cm, y=1cm]
    \node[storynode] (u) at (0,1.4) {$u$};
    \node[storynode] (v) at (3.6,0.2) {$v$};
    \node[storygoal] (t) at (1.8,-1.8) {$t$};
    \draw[storyedge, -{Stealth[length=2mm]}] (u) -- node[above right=1pt,
      font=\small] {$w(u,v)$} (v);
    \draw[alggreedy, line width=0.8pt, dashed, -{Stealth[length=2mm]}] (v) --
      node[right=3pt, font=\small] {$h(v)$} (t);
    \draw[alggreedy, line width=0.8pt, dashed, -{Stealth[length=2mm]}] (u) --
      node[left=3pt, font=\small] {$h(u)$} (t);
  \end{tikzpicture}
  \caption{Consistency as a triangle inequality: $h(u) \le w(u,v) + h(v)$ for
    every edge $(u,v)$.}\label{fig:consistency}
\end{figure}

Consistency is the stronger condition.

\begin{lemma}\label{lem:cons-adm}
Every consistent heuristic is admissible.
\end{lemma}

\begin{proof}
Let $h$ be a consistent heuristic. By Definition~\ref{def:heuristic}, $h$
takes nonnegative values, so it remains to show that $h(v) \le h^{*}(v)$ for
every vertex $v$. If $t$ is not reachable from $v$, then $h^{*}(v) = \infty$
and there is nothing to prove. Otherwise let $(v = u_0, u_1, \dots, u_m = t)$ be a shortest path from
$v$ to $t$. We show by induction on $m - i$ that
$h(u_i) \le \delta(u_i, t)$. For $i = m$ we have $h(t) = 0 = \delta(t,t)$. For
$i < m$, Inequality~\eqref{eq:consistent} and the induction hypothesis give
\[
  h(u_i) \le w(u_i, u_{i+1}) + h(u_{i+1})
         \le w(u_i, u_{i+1}) + \delta(u_{i+1}, t) = \delta(u_i, t),
\]
where the last equality holds because every part of a shortest path is itself a
shortest path. For $i = 0$ this is $h(v) \le h^{*}(v)$.
\end{proof}

The converse is false; Example~\ref{ex:inconsistent} gives an admissible
heuristic that is not consistent.

\subsection{Distance heuristics}\label{sec:distance-h}

On maps and grids the usual heuristics are distances between the coordinates
of a vertex $v = (x_v, y_v)$ and of the destination $t = (x_t, y_t)$.

\begin{definition}[Manhattan and Euclidean distance]\label{def:distances}
The \emph{Manhattan distance} and the \emph{Euclidean distance} from $v$ to $t$
are
\begin{align}
  h_{M}(v) &= |x_v - x_t| + |y_v - y_t|, \label{eq:manhattan}\\
  h_{E}(v) &= \sqrt{(x_v - x_t)^2 + (y_v - y_t)^2}, \label{eq:euclidean}
\end{align}
respectively. In three dimensions each formula gains a term for the third
coordinate.
\end{definition}

The Manhattan distance follows the axes, like a car on a street grid; the
Euclidean distance is the straight line. For a displacement of $(4, 3, 2)$ they
are 9 and $\sqrt{29} \approx 5.39$ (Figure~\ref{fig:euclid-box}). As surfaces
around the destination they form a pyramid and a cone
(Figure~\ref{fig:cone-pyramid}), and the pyramid never lies below the cone.

\begin{figure}[htbp]
  \centering
  \begin{tikzpicture}[x={(1.22cm,0cm)}, y={(0.66cm,0.48cm)},
                    z={(0cm,1.22cm)}, line join=round, line cap=round]
  \foreach \i in {0,...,4}{\draw[gridgray, line width=0.4pt] (\i,0,0) -- (\i,3,0);}
  \foreach \j in {0,...,3}{\draw[gridgray, line width=0.4pt] (0,\j,0) -- (4,\j,0);}
  \draw[gridgray, dashed, line width=0.6pt] (0,3,0) -- (0,3,2);
  \fill[rbluelt, opacity=0.8] (0,0,0) -- (4,3,0) -- (4,3,2) -- cycle;
  \draw[rbluemd, dashed, line width=1pt] (0,0,0) -- (4,3,0)
    node[pos=0.6, below right=-1pt, tag, text=rbluemd] {$5$};
  \fill[rbluelt, opacity=0.35] (0,0,2) -- (4,0,2) -- (4,3,2) -- (0,3,2) -- cycle;
  \draw[rbluedk!70, line width=0.8pt]
    (0,0,0) -- (4,0,0) -- (4,0,2) -- (0,0,2) -- cycle
    (4,0,0) -- (4,3,0) -- (4,3,2) -- (4,0,2)
    (0,0,2) -- (0,3,2) -- (4,3,2);
  \draw[cbOrange, line width=2.2pt] (0,0,0) -- (4,0,0) -- (4,3,0) -- (4,3,2);
  \node[tag, text=cbOrange, below=3pt] at (2,0,0) {$\Delta x = 4$};
  \node[tag, text=cbOrange, below right=0pt] at (4,1.5,0) {$\Delta y = 3$};
  \node[tagbox, text=cbOrange, left=3pt] at (4,3,1) {$\Delta z = 2$};
  \draw[rblue, line width=2.2pt] (0,0,0) -- (4,3,2);
  \node[tagbox, text=rblue, above left=0pt] at (2,1.5,1) {$\sqrt{29}\approx 5.39$};
  \node[sball] at (0,0,0) {};
  \node[tag, text=gstart, left=6pt] at (0,0,0) {\textbf{A}};
  \node[gball] at (4,3,2) {};
  \node[tag, text=ggoal, above=6pt] at (4,3,2) {\textbf{B}};
\end{tikzpicture}

  \caption{The Manhattan distance (orange, along the edges of the box) and the
    Euclidean distance (blue, straight through the box) from $A$ to
    $B$.}\label{fig:euclid-box}
\end{figure}

\begin{figure}[htbp]
  \centering
  \begin{tabular}{@{}c@{\hspace{1.5em}}c@{}}
  {\footnotesize\textcolor{rblue}{\textbf{Euclidean}} $h=\sqrt{x^2+y^2}$: a cone} &
  {\footnotesize\textcolor{cbOrange}{\textbf{Manhattan}} $h=|x|+|y|$: a pyramid}\\[2pt]
  \begin{tikzpicture}
    \begin{axis}[land3d, land view={-30}{42}{0.60},
                 colormap name=hblue, point meta min=0, point meta max=3.2]
      \addplot3[tilesurf, domain=0:4, y domain=0:360, samples=9, samples y=41]
        ({x*cos(y)}, {x*sin(y)}, {\kE*x});
      \node[gball] at (0,0,0) {};
      \node[tag, text=ggoal, below=6pt] at (0,0,0) {$t$};
      \node[vball] at (\pEx,\pEy,{\kE*veclen(\pEx,\pEy)}) {};
      \node[tagbox, text=rblue, anchor=west, right=3pt]
        at (\pEx,\pEy,{\kE*veclen(\pEx,\pEy)}) {$v$: $h=\sqrt5\approx2.24$};
    \end{axis}
  \end{tikzpicture}
  &
  \begin{tikzpicture}
    \begin{axis}[land3d, land view={-30}{42}{0.60},
                 colormap name=horange, point meta min=0, point meta max=3.2]
      \addplot3[tilesurf, domain=0:4, y domain=0:360, samples=9, samples y=41]
        ({x*cos(y)/(abs(cos(y))+abs(sin(y)))},
         {x*sin(y)/(abs(cos(y))+abs(sin(y)))}, {\kE*x});
      \node[gball] at (0,0,0) {};
      \node[tag, text=ggoal, below=6pt] at (0,0,0) {$t$};
      \node[vball] at (\pEx,\pEy,{\kE*(abs(\pEx)+abs(\pEy))}) {};
      \node[tagbox, text=cbOrange, anchor=west, right=3pt]
        at (\pEx,\pEy,{\kE*(abs(\pEx)+abs(\pEy))}) {$v$: $h=3$};
    \end{axis}
  \end{tikzpicture}\\
  {\scriptsize level sets are circles} & {\scriptsize level sets are diamonds}
\end{tabular}

  \caption{The two distance heuristics as surfaces around the destination $t$. At the
    point $v$ with offset $(2,1)$ from $t$ the pyramid is higher: 3 against
    $\sqrt{5} \approx 2.24$.}\label{fig:cone-pyramid}
\end{figure}

\begin{proposition}\label{prop:manh-eucl}
For any two points $v$ and $t$ of the plane, $h_{E}(v) \le h_{M}(v)$, with
equality if and only if $v$ and $t$ agree in at least one coordinate.
\end{proposition}

\begin{proof}
Let $a = |x_v - x_t|$ and $b = |y_v - y_t|$. Both distances are nonnegative, so
it suffices to compare their squares. We have
$h_{M}(v)^2 = (a+b)^2 = a^2 + b^2 + 2ab = h_{E}(v)^2 + 2ab$. Since
$2ab \ge 0$, it follows that $h_{E}(v) \le h_{M}(v)$, with equality exactly when
$ab = 0$, that is, when $a = 0$ or $b = 0$.
\end{proof}

Table~\ref{tab:which-h} shows where each distance is admissible. On a unit-cost
grid with four neighbors per cell, the Manhattan distance equals the true number
of moves when no obstacle intervenes. With straight-line edge weights, no chain
of edges is shorter than the straight line, so the straight-line distance is
admissible, while the larger Manhattan distance may overestimate.

\begin{table}[htbp]
  \centering
  \caption{When a distance heuristic never overestimates.}\label{tab:which-h}
  \renewcommand{\arraystretch}{1.25}%
  \begin{tabular}{@{}P{5.4cm}P{3.6cm}P{4.6cm}@{}}
    \toprule
    Graph & Admissible heuristic & Reason\\
    \midrule
    Grid with four neighbors per cell, every move of cost 1
      (Section~\ref{sec:grid})
      & Manhattan distance & every move changes one coordinate by 1\\
    Edge weights equal to straight-line distances
      (Section~\ref{sec:mapgraph})
      & Euclidean or great-circle distance & no chain of edges is shorter
        than the direct line\\
    Edge weights unrelated to the coordinates
      (Section~\ref{sec:overestimate})
      & neither, in general & a distance can exceed the true cost\\
    \bottomrule
  \end{tabular}
\end{table}

In both cases the admissible heuristic is also consistent: neighboring grid
cells differ by exactly 1 in Manhattan distance, and a straight line is never
longer than a detour through a neighbor.

\subsection{How a heuristic guides the search}\label{sec:landscape}

Lifting every vertex to height $h(v)$ turns the heuristic into a funnel around
the destination (Figure~\ref{fig:landscape-a}). Table~\ref{tab:slope} shows how
the key changes for a unit move on a grid with the Euclidean heuristic
(Figure~\ref{fig:landscape-b}).

\begin{figure}[htbp]
  \centering
  \begin{subfigure}[b]{0.50\textwidth}\centering
    \begin{tikzpicture}
  \begin{axis}[land3d, land view={-32}{34}{0.42}]
    \floorA
    \addplot3[tilesurf, colormap name=hblue, samples=13, samples y=13,
              domain=0:12, y domain=0:12]
      {\kA*veclen(x-\gAx, y-\gAy)};
    \draw[neutralg, line width=0.6pt, {Bar[width=1.6mm]}-{Bar[width=1.6mm]}]
      (12,0,0) -- \liftA{12}{0}
      node[tagbox, text=neutralg, midway, left=2pt] {$h(v)$};
    \draw[edgpath, line cap=round, line join=round]
      \liftA{2}{10} -- \liftA{3}{9} -- \liftA{4}{8} -- \liftA{5}{7}
      -- \liftA{6}{6} -- \liftA{7}{5} -- \liftA{7}{4} -- \liftA{7}{3};
    \draw[move3d,cbRed] \liftA{2}{10} -- \liftA{1}{11};
    \node[tagbox, text=cbRed, above=3pt] at \liftA{1}{11} {uphill: $\Delta f=+2$};
    \node[tagbox, text=cbGreen, left=5pt] at \liftA{4}{8} {downhill: $\Delta f=0$};
    \node[sball] at \liftA{2}{10} {};
    \node[tag, text=gstart, below left=3pt] at \liftA{2}{10} {$s$};
    \node[gball] at \liftA{7}{3} {};
    \node[tag, text=ggoal, below=6pt] at \liftA{7}{3} {$t$, \ $h=0$};
  \end{axis}
\end{tikzpicture}

    \caption{A $12 \times 12$ map lifted by the Euclidean heuristic. The path
      from $s$ runs downhill to $t$.}\label{fig:landscape-a}
  \end{subfigure}\hfill
  \begin{subfigure}[b]{0.46\textwidth}\centering
    \begin{tikzpicture}
  \begin{axis}[land3d, land view={-24}{30}{0.62}]
    \addplot3[tilesurf, colormap name=hblue, samples=10, samples y=7,
              domain=0:9, y domain=-3:3] {\kB*veclen(x, y)};
    \node[gball] at \liftB{0}{0} {};
    \node[tag, text=ggoal, left=6pt] at \liftB{0}{0} {$t$};
    \draw[cbGreen, dashed, line width=1pt] \liftB{5}{0} -- \liftB{0.5}{0};
    \draw[move3d,cbGreen]  \liftB{6}{0} -- \liftB{5}{0};
    \draw[move3d,cbRed]    \liftB{6}{0} -- \liftB{7}{0};
    \draw[move3d,neutralg] \liftB{6}{0} -- \liftB{6}{1};
    \draw[move3d,neutralg] \liftB{6}{0} -- \liftB{6}{-1};
    \node[vball] at \liftB{6}{0} {};
    \node[tag, above left=4pt] at \liftB{6}{0} {$v$};
    \node[tagbox, text=cbGreen, anchor=north] at \liftB{4.7}{-0.35} {$\Delta f = 0$};
    \node[tagbox, text=cbRed, anchor=west] at \liftB{7.15}{0} {$\Delta f = +2$};
    \node[tagbox, text=neutralg, anchor=west] at \liftB{6.2}{-1.25} {$\Delta f \approx +1$};
  \end{axis}
\end{tikzpicture}

    \caption{Four unit moves out of a vertex $v$ at distance 6 from
      $t$.}\label{fig:landscape-b}
  \end{subfigure}
  \caption{The heuristic as a landscape.}\label{fig:landscape}
\end{figure}

\begin{table}[htbp]
  \centering
  \caption{Change of the key for a unit move out of a vertex at distance 6 from
    the destination, with the Euclidean heuristic.}\label{tab:slope}
  \begin{tabular}{@{}lccc@{}}
    \toprule
    Move & $\Delta g$ & $\Delta h$ & $\Delta f$\\
    \midrule
    toward the destination    & $+1$ & $-1$ & \textcolor{algastar}{$0$}\\
    away from the destination & $+1$ & $+1$ & \textcolor{cbRed}{$+2$}\\
    sideways           & $+1$ & $\sqrt{37} - 6 \approx +0.08$ & $\approx +1.08$\\
    \bottomrule
  \end{tabular}
\end{table}

A move toward the destination leaves the key unchanged, while a move away raises
it by 2. Such moves are not forbidden, only postponed, so A* expands vertices in
the direction of the destination first.

\subsection{Misleading heuristics}\label{sec:bad-h}

Table~\ref{tab:bad-h} lists three ways in which a heuristic can mislead A*.
Only an overestimating heuristic can cost the shortest path; the other two only
waste work.

\begin{table}[htbp]
  \centering
  \caption{Three ways in which a heuristic can mislead A*.}\label{tab:bad-h}
  \renewcommand{\arraystretch}{1.25}%
  \begin{tabular}{@{}P{2.6cm}P{3.4cm}P{2.1cm}P{5.6cm}@{}}
    \toprule
    Kind & Property & Shortest path & Effect\\
    \midrule
    \textcolor{cbRed}{Overestimating} & not admissible & not guaranteed
      & A* behaves like greedy best-first search; on the running example the
        returned cost rises from 8 to 10 (Section~\ref{sec:overestimate})\\
    Uninformative & admissible, nearly constant & guaranteed
      & $h(v) = 1$ for every $v \neq t$ (admissible if every edge costs at
        least 1) adds the same constant to almost every key, so A* expands
        almost the same vertices as Dijkstra's algorithm\\
    \textcolor{alggreedy}{Blind to obstacles} & admissible, ignores walls
      & guaranteed & A* expands dead ends that appear close to the destination, such
        as the inside of the U-shaped wall of Section~\ref{sec:grid}\\
    \bottomrule
  \end{tabular}
\end{table}

\subsubsection{An overestimating heuristic}\label{sec:overestimate}

We place the cities of the running example on a grid at $\cty{D} = (0,4)$, $\cty{N} = (4,6)$,
$\cty{B} = (8,4)$, $\cty{H} = (12,6)$, $\cty{S} = (16,4)$, $\cty{G} = (2,1)$,
$\cty{K} = (8,8)$ and $\cty{C} = (13,1)$ (Figure~\ref{fig:story-m}), and use the Manhattan distance $h_{M}$ to \cty{S} as the heuristic. The grid
units are unrelated to the edge weights, and $h_{M}$ overestimates in six of the
eight cities (Table~\ref{tab:story-hm}).

\begin{table}[htbp]
  \centering
  \caption{Two heuristics for the running example and the true remaining cost.
    Red values overestimate.}\label{tab:story-hm}
  \begin{tabular}{@{}lcccccccc@{}}
    \toprule
    City & \cty{D} & \cty{G} & \cty{N} & \cty{K} & \cty{B} & \cty{C} & \cty{H} & \cty{S}\\
    \midrule
    Heuristic $h$ (Table~\ref{tab:story-h}) & 7 & 8 & 5 & 4 & 4 & 7 & 2 & 0\\
    Manhattan distance $h_{M}$ & \textcolor{cbRed}{16} & \textcolor{cbRed}{17}
      & \textcolor{cbRed}{14} & \textcolor{cbRed}{12} & \textcolor{cbRed}{8}
      & 6 & \textcolor{cbRed}{6} & 0\\
    True cost $h^{*}$     & 8 & 9 & 6 & 7 & 5 & 8 & 3 & 0\\
    \bottomrule
  \end{tabular}
\end{table}

With $h_{M}$, A* returns $(\cty{D}, \cty{B}, \cty{H}, \cty{S})$ of cost 10
(Table~\ref{tab:trace-m}). The cause is $h_{M}(\cty{N}) = 14$, more than twice
$h^{*}(\cty{N}) = 6$: after step 1, \cty{N} has key 16 and \cty{B} key 13, so
the search reaches \cty{S} with $g = 10$ before expanding \cty{N}. With
$h(\cty{N}) = 5$, the key of \cty{N} would be $7 < 10$. When $h$ dominates $g$,
A* behaves like greedy best-first search and returns the same path.

\begin{table}[htbp]
  \centering
  \caption{A* on the running example with the Manhattan distance $h_{M}$, which
    overestimates.}\label{tab:trace-m}
  \small
  \begin{tabular}{@{}c>{\raggedright}P{4.8cm}cc>{\raggedright\arraybackslash}P{3.4cm}@{}}
    \toprule
    Step & $Q$ before extraction (vertex: key) & Extracted & $g + h_{M} = f$
      & Updates\\
    \midrule
    \inputrows{figures/trace_manhattan.tex}
    \bottomrule
  \end{tabular}
\end{table}

\begin{figure}[htbp]
  \centering
  \resizebox{0.62\textwidth}{!}{\begin{tikzpicture}[x=1cm, y=1cm]
  \fill[black!10] (-0.529,-0.393) -- (7.052,-1.657) -- (8.132,1.475) -- (0.551,2.739) -- cycle;
  \draw[draw=rblue!14, fill=rbluelt!35, line width=0.6pt] (-0.529,-0.243) -- (7.052,-1.507) -- (8.132,1.625) -- (0.551,2.889) -- cycle;
  \draw[black!19, line width=0.3pt] (0.000,0.000) -- (0.864,2.506);
  \draw[black!10, line width=0.3pt] (0.421,-0.070) -- (1.285,2.435);
  \draw[black!10, line width=0.3pt] (0.842,-0.140) -- (1.706,2.365);
  \draw[black!10, line width=0.3pt] (1.264,-0.211) -- (2.128,2.295);
  \draw[black!19, line width=0.3pt] (1.685,-0.281) -- (2.549,2.225);
  \draw[black!10, line width=0.3pt] (2.106,-0.351) -- (2.970,2.155);
  \draw[black!10, line width=0.3pt] (2.527,-0.421) -- (3.391,2.084);
  \draw[black!10, line width=0.3pt] (2.948,-0.491) -- (3.812,2.014);
  \draw[black!19, line width=0.3pt] (3.370,-0.562) -- (4.234,1.944);
  \draw[black!10, line width=0.3pt] (3.791,-0.632) -- (4.655,1.874);
  \draw[black!10, line width=0.3pt] (4.212,-0.702) -- (5.076,1.804);
  \draw[black!10, line width=0.3pt] (4.633,-0.772) -- (5.497,1.733);
  \draw[black!19, line width=0.3pt] (5.054,-0.842) -- (5.918,1.663);
  \draw[black!10, line width=0.3pt] (5.476,-0.913) -- (6.340,1.593);
  \draw[black!10, line width=0.3pt] (5.897,-0.983) -- (6.761,1.523);
  \draw[black!10, line width=0.3pt] (6.318,-1.053) -- (7.182,1.453);
  \draw[black!19, line width=0.3pt] (6.739,-1.123) -- (7.603,1.382);
  \draw[black!19, line width=0.3pt] (0.000,0.000) -- (6.739,-1.123);
  \draw[black!10, line width=0.3pt] (0.108,0.313) -- (6.847,-0.810);
  \draw[black!10, line width=0.3pt] (0.216,0.626) -- (6.955,-0.497);
  \draw[black!10, line width=0.3pt] (0.324,0.940) -- (7.063,-0.184);
  \draw[black!19, line width=0.3pt] (0.432,1.253) -- (7.171,0.130);
  \draw[black!10, line width=0.3pt] (0.540,1.566) -- (7.279,0.443);
  \draw[black!10, line width=0.3pt] (0.648,1.879) -- (7.387,0.756);
  \draw[black!10, line width=0.3pt] (0.756,2.192) -- (7.495,1.069);
  \draw[black!19, line width=0.3pt] (0.864,2.506) -- (7.603,1.382);
  \coordinate (D) at (0.432,1.253);
  \coordinate (N) at (2.333,1.598);
  \coordinate (B) at (3.802,0.691);
  \coordinate (H) at (5.702,1.037);
  \coordinate (S) at (7.171,0.130);
  \coordinate (G) at (0.950,0.173);
  \coordinate (K) at (4.234,1.944);
  \coordinate (C) at (5.584,-0.599);
  \draw[draw=rblue!55, line width=1pt] (D) -- (N);
  \draw[draw=alggreedy, line width=2pt] (D) -- (B);
  \draw[draw=rblue!55, line width=1pt] (N) -- (B);
  \draw[draw=alggreedy, line width=2pt] (B) -- (H);
  \draw[draw=alggreedy, line width=2pt] (H) -- (S);
  \draw[draw=rblue!55, line width=1pt] (D) -- (G);
  \draw[draw=rblue!55, line width=1pt] (N) -- (K);
  \draw[draw=rblue!55, line width=1pt] (K) -- (H);
  \draw[draw=rblue!55, line width=1pt] (B) -- (C);
  \draw[draw=rblue!55, line width=1pt] (C) -- (S);
  \node[isoweight] at ($(D)!0.5!(N)$) {2};
  \node[isoweight] at ($(D)!0.5!(B)$) {5};
  \node[isoweight] at ($(N)!0.5!(B)$) {1};
  \node[isoweight] at ($(B)!0.5!(H)$) {2};
  \node[isoweight] at ($(H)!0.5!(S)$) {3};
  \node[isoweight] at ($(D)!0.5!(G)$) {1};
  \node[isoweight] at ($(N)!0.5!(K)$) {3};
  \node[isoweight] at ($(K)!0.5!(H)$) {4};
  \node[isoweight] at ($(B)!0.5!(C)$) {3};
  \node[isoweight] at ($(C)!0.5!(S)$) {10};
  \isodisc{D}{alggreedy!15}{algastar}{D}{16}
  \isodisc{N}{white}{rblue!70}{N}{14}
  \isodisc{B}{alggreedy!15}{alggreedy}{B}{8}
  \isodisc{H}{alggreedy!15}{alggreedy}{H}{6}
  \isodisc{S}{alggreedy!15}{cbRed}{S}{0}
  \isodisc{G}{white}{rblue!70}{G}{17}
  \isodisc{K}{white}{rblue!70}{K}{12}
  \isodisc{C}{white}{rblue!70}{C}{6}
\end{tikzpicture}
}
  \caption{The running example on its grid, with the Manhattan distance $h_{M}$
    in each disc. A* now returns the orange path, which is not the shortest
    one.}\label{fig:story-m}
\end{figure}

Figure~\ref{fig:estimates} compares three estimates for \cty{N}; only the one
above the true cost leads A* astray.

\begin{figure}[htbp]
  \centering
  \begin{tikzpicture}[x=0.8cm, y=1cm]
    \fill[algastar!15] (0,-0.12) rectangle (6,0.12);
    \fill[cbRed!10] (6,-0.12) rectangle (15.2,0.12);
    \draw[inkgray, line width=0.6pt, -{Stealth[length=2mm]}] (0,0) -- (15.6,0);
    \foreach \x in {0,2,...,14}
      {\draw[inkgray] (\x,-0.08) -- (\x,0.08);
       \node[font=\scriptsize, text=inkgray, below=3pt] at (\x,-0.08) {\x};}
    \draw[algastar, line width=1pt] (6,-0.25) -- (6,0.25);
    \node[font=\small, text=algastar, below=16pt] at (6,0) {$h^{*}(\cty{N}) = 6$};
    \foreach \x/\c/\lab/\note in {0/inkgray/{$h = 0$}/{no information},
                                  5/algastar/{$h(\cty{N}) = 5$}/{admissible},
                                  14/cbRed/{$h_M(\cty{N}) = 14$}/{overestimate}}
      {\draw[\c, line width=0.9pt, -{Stealth[length=2mm]}] (\x,0.85) -- (\x,0.18);
       \fill[\c] (\x,0) circle (2.2pt);
       \node[font=\small, text=\c, above] at (\x,0.85) {\lab};
       \node[font=\scriptsize, text=\c, above=13pt] at (\x,0.85) {\note};}
    \node[font=\footnotesize, text=algastar] at (3,-1.25) {admissible range:
      $0 \le h(\cty{N}) \le h^{*}(\cty{N})$};
    \node[font=\footnotesize, text=cbRed] at (10.6,-1.25) {overestimates};
  \end{tikzpicture}
  \caption{Three estimates of the remaining cost from \cty{N}, compared with the
    true cost $h^{*}(\cty{N}) = 6$.}\label{fig:estimates}
\end{figure}
